\documentclass[11pt]{article}
\input{style-pc}
\usepackage{array}

\begin{document}

\entetesujet{Sujet bac 2016 -- Physique-Chimie -- Série D}

% ============================ CHIMIE ============================
\matiere{CHIMIE}{8 points}

\partie{A : vérification des connaissances}

\noindent\textbf{1. Texte à trous.} Complète la phrase ci-après par les mots suivants : \textit{neutrons, l'atome, nucléons, protons}.

\textit{Le noyau $\cdots\cdots$ est constitué des $\cdots\cdots$ et des $\cdots\cdots$ appelés $\cdots\cdots$}

\medskip
\noindent\textbf{2. Appariement.} Relie un élément-question de la colonne A à un élément-réponse de la colonne B. Exemple : $A_7 = b_8$.

\begin{center}\renewcommand{\arraystretch}{1.6}
\begin{tabular}{|p{0.4\linewidth}|p{0.5\linewidth}|}
\hline
\textbf{Colonne A} & \textbf{Colonne B} \\
\hline
$a_1$ : solution basique & $b_1$ : contient plus d'ions hydronium que d'ions hydroxyde \\
\hline
$a_2$ : solution acide & $b_2$ : isotopes \\
\hline
$a_3$ : le produit ionique de l'eau & $b_3$ : contient plus d'ions hydroxyde que d'ions hydronium \\
\hline
$a_4$ : nucléides de même $Z$ mais de $A$ différents & $b_4$ : aqueuse diluée \\
\hline
\end{tabular}
\end{center}

\noindent\textbf{3. Question à réponse construite.} La radioactivité $\gamma$ accompagne généralement les radioactivités $\alpha$ et $\beta$. Explique l'émission du rayonnement $\gamma$.

\partie{B : application des connaissances}
On a préparé une solution d'acide méthanoïque ($\ch{HCOOH}$) de concentration inconnue $C_0$. Le pH de cette solution est égal à $2{,}7$.

\begin{enumerate}
  \item Écris l'équation-bilan de la réaction de l'acide méthanoïque avec l'eau.
  \item Fais l'inventaire de toutes les espèces chimiques présentes dans la solution.
  \item Sachant que le pKa du couple acide-base associé à l'acide méthanoïque est $3{,}8$, calcule les concentrations molaires de toutes les espèces chimiques.
  \item Déduis la concentration molaire $C_0$ de la solution d'acide méthanoïque.
\end{enumerate}
\noindent On donne $K_e = 10^{-14}$.

% ============================ PHYSIQUE ============================
\matiere{PHYSIQUE}{12 points}

\partie{A : vérification des connaissances}

\noindent\textbf{1. Question à réponse courte.} Définis une onde transversale.

\medskip
\noindent\textbf{2. Réarrangement.} Ordonne la phrase suivante qui est écrite en désordre.

\textit{Un oscillateur / est une fonction sinusoïdale du temps / est / dont / un oscillateur / harmonique / l'équation horaire du mouvement.}

\medskip
\noindent\textbf{3. Questions à alternative vrai ou faux.} Réponds par vrai ou faux aux affirmations suivantes :
\begin{enumerate}[label=\alph*.]
  \item Pour des oscillations de faible amplitude, un pendule pesant est un oscillateur harmonique de translation.
  \item La longueur d'onde est la distance parcourue par une onde en une période.
\end{enumerate}

\medskip
\noindent\textbf{4. Schéma à compléter.} Dans le repère ci-dessous, représente les vecteurs de Fresnel $\vv{OA_1}$ et $\vv{OA_2}$ associés respectivement aux fonctions sinusoïdales : $y_1 = 2\sin\!\left(100\pi t + \dfrac{\pi}{2}\right)$ (cm) et $y_2 = 3\sin(50\pi t + \pi)$ (cm).

\begin{center}
\begin{tikzpicture}[>=Stealth,scale=1]
  \draw[->] (-2.6,0)--(2.6,0) node[right]{$x$(cm)};
  \draw[->] (0,-2.6)--(0,2.6) node[above]{$y$(cm)};
  \node[above left] at (-2.6,0){$x'$};
  \node[below right] at (0,-2.6){$y'$};
  \draw[->] (0,0)--(0.7,0) node[below]{$\vi$};
  \draw[->] (0,0)--(0,0.7) node[left]{$\vj$};
  \node[below left] at (0,0){$o$};
\end{tikzpicture}
\end{center}

\partie{B : application des connaissances}
Une corde de longueur $L = 1{,}2$ m et de masse $m = 40$ g soutient à son extrémité $C$ un solide de masse $M$. La partie horizontale de longueur $l = AB = 0{,}9$ m est le siège d'un phénomène d'ondes stationnaires. Un vibreur impose à l'extrémité $A$ un mouvement sinusoïdal transversal de fréquence $N = 50$ Hz. L'extrémité $C$ de la corde porte un solide de masse $M$.

\begin{center}
\begin{tikzpicture}[>=Stealth,scale=1]
  % vibreur
  \draw[pattern=north east lines] (-0.9,-0.12) rectangle (0,0.12);
  \node[left] at (-0.9,0){vibreur};
  \node[above] at (0.4,0.1){$A$};
  % corde horizontale
  \draw (0,0)--(5.5,0);
  % poulie
  \draw (6,0) circle (0.5);
  \fill (6,0) circle (1pt);
  \node[above] at (6,0.55){$B$};
  % corde verticale vers C
  \draw (6.5,0)--(6.5,-2.3);
  % masse
  \draw[fill=gray!30] (6.32,-2.6) rectangle (6.68,-2.3);
  \node[right] at (6.55,-2.1){$C$};
  \node[right] at (6.7,-2.45){$(M)$};
\end{tikzpicture}
\end{center}

\begin{enumerate}
  \item On observe 6 fuseaux sur la partie $AB$ de la corde. Calcule :
  \begin{enumerate}[label=\alph*.]
    \item La longueur d'onde de la vibration.
    \item La célérité de propagation.
  \end{enumerate}
  \item Détermine la valeur de la masse $M$.
\end{enumerate}
\noindent On donne $g = 10$ m/s$^2$.

\partie{C : résolution d'un problème}
On se propose de déterminer la vitesse d'un électron à la sortie des plaques d'un condensateur où règne un champ électrostatique uniforme $\vec{E}$. Pour cela, on considère un faisceau homocinétique d'électrons qui pénètre au point $O$ avec une vitesse initiale $V_0 = 10^7$ m/s, dans le champ électrostatique uniforme $\vec{E}$ compris entre deux plaques métalliques parallèles et horizontales $A$ et $B$ distantes de $d = 15$ cm.

\begin{center}
\begin{tikzpicture}[>=Stealth,scale=1]
  \draw[->] (0,-0.2)--(0,2) node[above]{$y$};
  \draw[->] (0,0)--(6.4,0) node[right]{$x$};
  \draw[thick] (0.3,1)--(4.3,1) node[right]{$A$};
  \draw[thick] (0.3,-1)--(4.3,-1) node[right]{$B$};
  \node[below left] at (0,0){$o$};
  \draw[->,thick] (0,0)--(1.2,0) node[below]{$\vec{v_0}$};
  \draw[->] (0,0)--(0.6,0); \draw[->] (0,0)--(0,0.6) node[left]{$\vj$}; \node[below] at (0.7,0){$\vi$};
  \fill (4,0.55) circle (1.5pt) node[above right]{$S$};
  \draw[dashed] (0,0)--(4,0.55);
  \draw[<->,dashed] (-0.3,-1)--(-0.3,1) node[midway,left]{$d$};
  \draw[<->,dashed] (0.3,-1.4)--(4.3,-1.4) node[midway,below]{$l$};
\end{tikzpicture}
\end{center}

\begin{enumerate}
  \item On établit entre ces plaques, de longueur $l = 20$ cm, une différence de potentiel $V_A - V_B = U_{AB} = +150$ V.
  \begin{enumerate}[label=\alph*.]
    \item Donne le signe des plaques $A$ et $B$ puis représente le vecteur champ $\vec{E}$ entre les plaques.
    \item Calcule l'intensité du vecteur champ $\vec{E}$.
  \end{enumerate}
  \item \begin{enumerate}[label=\alph*.]
    \item Établis les équations horaires du mouvement de l'électron dans le repère $(O,\vi,\vj)$.
    \item Détermine l'équation cartésienne de la trajectoire.
  \end{enumerate}
  \item Trouve les composantes du vecteur vitesse $\vec{v_S}$ à la sortie $S$ du champ.
  \item Déduis alors la norme du vecteur vitesse $\vec{v_S}$.
\end{enumerate}
\noindent Données : $q = e = -1{,}6\cdot10^{-19}$ C ; $m = 9{,}1\cdot10^{-31}$ kg.

\end{document}
